\section{Motivation}
\label{section:41-GAI-Motivation}

% \begin{figure}[h!]
    \begin{center}
        \includesvg[width=0.66\textwidth]{Figures/SVGs/xkcd}
        \caption{XKCD `Standards'. Source: \url{https://xkcd.com/927/}}
        \label{figure:xkcd}
    \end{center}
\end{figure}

% \nrscomment{make this active- in reviewing x we conclude y}
In reviewing prior work in competitive benchmarking and workload analysis (section~\ref{section:31-literature-benchmarking}) we conclude that purely competitive benchmarks from the Input/Output paradigm lack descriptive power and that empirical benchmarks from the workload analyses paradigm are too invasive for general-purpose use. 
\nrscomment{reword: "To achieve a balance between ease-of-use and descriptive power, we propose a hybrid approach inspired by constrained benchmarks that includes external levers to enable a detailed black box measurement approach."}
Utility must balance explainability, and our hybrid approach borrows ideas from constrained benchmarks, creating external levers to facilitate a more detailed `black-box' approach to understanding the internal state of the system.

The review of \gai benchmarking (section~\ref{section:32-literature-benchmarking-gai}) \nrscomment{also} concludes that the existing \gai task, dataset, and metric coverage is narrow, neglects dynamic processing of heterogeneous graphs, and fails to appropriately evaluate approximate and incremental solutions. 
\nrscomment{missing the link between the above issues and the below solution to them. Why does a problem-centric instead of system-centric approach fix dataset/metric/task/dynamic graph/heterogeneous graph coverage? By the end of this para you explain covering the breadth of datasets/metrics/tasks in the problem domain, so thats that link (though it needs a better transition so by this point I get there will be a connection). But how does problem-centric also help tackle the dynamic/heterogeneous/incremental issue? By listing those above you tee up addressing them and then don't.}

\nrscomment{transition: "To address these issues, we follow a problem-centric design, rather than the typical system-centric approach. For each relevant \gai problem category (see table...), we derive relevant tasks from a study of IandS domain reqs (See...) and define benchmark tasks, datasets, and metrics that cover the breadth of \gai implementations and environments." ......... (i had been expecting something problem-related to go here, since the approach is problem-centric and were talking about reviwing the problem reqs to define tasks/metrics/etc, but you have \gai implementations and evironments..why? what does it mean to have tasks/datasets/metrics that cover the breadth of \gai implementations and environments? isnt your whole thing that the tasks shouldn't presuppose an implementation or environment?) led approach, we take a problem-centric approach.
We focus on the \gai problem categories (see table~\ref{table:taxonomy}), derive relevant tasks within those categories from a study of I\&S domain requirements (see table:~\ref{table:problem-survey}), and design a benchmark framework that defines the tasks, datasets and metrics for evaluating across the breadth of \gai implementations and environments.

\nrscomment{this sentence has too many clauses. "Our flexible, problem-centric approach provides a unified framework for varying benchmark components like algorithm, implementation, and hardware while measuring corresponding changes in performance.}
The problem-centric approach offers the flexibility to vary any benchmark component from the \textit{algorithm} or \textit{implementation} to the underlying \textit{hardware} within a unified framework, while monitoring the changes in efficiency and effectiveness.

% \nrscomment{i think you can cut this entire chunk. much of it is about stuff you already cut from the previous section and the rest is redundant to what you just said}
% The review of \gai bottlenecks (section~\ref{section:33-literature-gai-bottlenecks}) concludes that the three methods to influence the effectiveness, efficiency or cost of a \gai system are \textit{algorithmic, implementation} and \textit{hardware} changes.
% First, the \textit{algorithms} used to analyze the graphs can be changed to approach a task differently, operate in a time-constrained setting, or approximate results. 
% Second, the underlying \textit{implementation} can be changed to better exploit existing hardware capabilities, use optimized libraries, and adjust the memory consumption (or access time) for a graph in memory. 
% The third major avenue for improvement is to modify the \textit{hardware} to increase computing resources, distribute computing or use specialized accelerators. 
% Existing Input/Output paradigm benchmarks focus tend to focus only on evaluating hardware changes, while hosted benchmarks only evaluate the algorithmic changes. 
% Workload analyses explore all three dimensions but are invasive and tend to ignore effectiveness metrics, in favor of a focus on efficiency. 
% A key goal of our benchmark is to describe the effectiveness, efficiency and cost trade-offs for choosing some combination of these acceleration methods across the breadth of the \gai tasks they need to undertake.

% \nrscomment{cut this, too colloquial}
% Whenever developing a new standard, the new approach risks crowding the already busy ecosystem, without meaningful improvement (see: figure~\ref{figure:xkcd}).
% The reasons why the framework we propose here is an important contribution and not just crowding existing work are: 


% \osullikomment{Move to introduction.}

% \nrscomment{Our \gai benchmarking methodology.... (id turn this into prose or figure out what this is a list of and call it that...goals? characteristics? guiding principles? properties?}

% \begin{enumerate}
%     \item Our framework defines a methodology, allowing it to extend to emerging \gai problems. 
%     \item Our methodology focuses on the \gai problems, rather than the systems. 
%     \item Our benchmark compares the trade-offs between the effectiveness, efficiency and cost of a \gai system to achieve a task. 
%     \item Our benchmark balances the generality of the Input/Output benchmarking paradigm with the insights of the empirical benchmarking paradigm.
% \end{enumerate}


% They can be \textit{competitive}, where candidate systems, algorithms, etc are \textit{`raced'} against each other, or \textit{empirical} benchmarks for the study of performance bottlenecks~\cite{Saavedra1996}.
% In the AI era, the public evaluation sets that are projected onto leaderboards are implicit \textit{competitive} benchmarks.
% These leaderboards typically report accuracy, harmonic mean, execution time or some other custom metric that measures the \textit{effectiveness} or \textit{efficency} of the AI tasks at hand, be that object detection, machine translation, common-sense reasoning or simple classification. 
% The mono-dimensionality of these leaderboards can conceal bias, fairness and generalizability of these models, with overfitting to problem datasets a known problem~\cite{Hooker1994, Hooker1995}.

% %
% \nrscomment{you told me there are 2 things in the world, then described one and not the other. talk about what an empirical benchmark is, even if just at a high level}
% %

% %
% \nrscomment{before you say this section is about the need for x, tell me theres a need for x because y. i think thats what this signpost ought to be, not all the stuff about leaderboards}

% \nrscomment{i think this section will do more than just outline the need for it, it will give a blueprint for addressing the need. say that}
% %
% This section outlines the need for a graph-AI benchmark, the methodology, design, implementation and results. 
